\subsection*{Training the open glyph writer}
A six-layer causal transformer groups four consecutive bitmap cells and uses a separate causal readout inside each group. It observes the pooled ink of 32-pixel cells. The density mixes a learned score over 5,545 distinct fitting bitmaps, a pointer to visible preceding ink, and termination. Repeated physical shapes share their encoding. A next output may include a previously unseen input bitmap through copying. We trained on 23 English and 20 Chinese works from GOLEM, 50,280 English and 47,496 Chinese synthetic instruction records from Alpaca, and 33,503 training examples from \href{https://winogrande.allenai.org/}{WinoGrande}. Five distinct works supplied the reported fiction tests. Cropping limits the evidence available to the model. Shared translations across the two instruction corpora remain possible. We screened eight-word anonymized passage overlap between training references and held families. The matched reference runs consumed identical batches. A byte-input control received byte values during fitting and inference. Each warm fit used 8,000 steps and each continued fit used 2,000. The deployed optical path took 182.9 seconds and 25.92 kJ on one RTX 4090 D, excluding data preparation, host energy, controls and evaluation.

For paired training let $d_1,d_2$ be the page-conditioned log odds of the same two written answers, with opposite correct choices. Both strict decisions are correct exactly when $d_1-d_2>|d_1+d_2|$. We checked the integer form in Lean and optimized a smooth penalty on the smaller signed margin after normalizing by answer length. Counts of self-generated answers included the emitted stop. For books, matching a single recorded continuation is a stringent criterion. The prompts in Fig.~\ref{fig:visual-discourse}a are the first held examples in a fixed hash order. A separate CPU replay directly read those PNG images and reproduced the emitted glyph stream.

For human question answering we retained the 2,117,313-parameter writer and added a 528,197-parameter reader and recurrent emitter. A visible source state denotes the position from which the last glyph was copied. Let $h_j$ be its posterior mass, $c_j$ the learned probability of advancing and $a_j$ a normalized start distribution. Define $v_{j+1}=h_j c_j$ at valid next cells and $r=1-\sum_j v_j$. The source emission weights are $v_j+r\beta a_j$, while the free emission has weight $r(1-\beta)$. A separate stop hazard precedes this mixture. Summing matching emission weights gives the target likelihood. Normalizing the compatible source weights gives the next posterior. The independent comparison sets $v_j=0$. Exhaustive enumeration of short paths checked the recursion. Both arms received 2,400 matched batches and 1,120,936 supervised answer cells each. The persistent addition used 118.0 seconds and 12.66 kJ on the same device.

We prepared 105,513 English \href{https://stanfordnlp.github.io/coqa/}{CoQA}, 9,981 simplified \href{https://aclanthology.org/D19-1600/}{CMRC} and 26,775 traditional \href{https://github.com/DRCKnowledgeTeam/DRCD}{DRCD} training examples. A page contains the current question, one preceding exchange for English dialogue when available, and the passage prefix fitting within 768 cells. No reference span guides the crop. Training targets are capped at 48 cells, without adding a stop to a censored target. We evaluated the first two eligible questions with distinct answers from each of 64 development passages per corpus in hash order. Free replies had a 96-cell cap. Scoring requires an emitted stop and a reference match after case folding, whitespace removal and boundary punctuation trimming. English dialogue used the single reference in the prepared release. The scores describe this paired sample.
